\section{Introdução}

\begin{frame}{O SPED e a Escrituração Fiscal Digital de ICMS e IPI}
\begin{itemize}
    \item \textbf{O SPED:} É um sistema que busca modernizar a interação entre contribuintes e o fisco através da informatização das obrigações acessórias.
    \item \textbf{O EFD ICMS IPI:} É de uso obrigatório para os contribuintes de ICMS e IPI e contém registros de apuração de impostos, documentos fiscais e outras informações de interesse.
    \item \textbf{A urgência legal:} Prazo de 5 anos para auditorias antes que seja extinto o crédito tributário.
\end{itemize}
\end{frame}

\begin{frame}{O Desafio Técnico Existente}
\textbf{Com isso:}
\begin{itemize}
    \item \textbf{Auditorias:} É de interesse do Fisco realizar auditorias sobre esses documentos em tempo hábil.
    \item \textbf{Empecilhos:} Arquivos em texto plano hierárquico e atualizações constantes de leiaute pelo Fisco.
    \item \textbf{O desafio gerado:} Alto custo de manutenção e quebra de sistemas com código engessado.
\end{itemize}
\end{frame}

\begin{frame}{A Nossa Proposta}
\textbf{Portanto, nossa equipe definiu como objetivo geral:}
\begin{itemize}
    \item Desenvolver um \textit{pipeline} ETL dinâmico e orientado a metadados para EFD ICMS IPI.
    \item Realizar a substituição de regras no código-fonte por contratos em formato \textbf{JSON}.
    \item Viabilizar a ingestão escalável dos documentos, com um processamento resiliente e genérico, junto com armazenamento relacional preciso com uso do \textit{PostgreSQL}.
\end{itemize}
\end{frame}